\documentclass[11pt]{article}
\usepackage[T1]{fontenc}
\usepackage{lmodern,amsmath,amssymb,amsthm,mathtools}
\usepackage[margin=1in]{geometry}
\usepackage{microtype,booktabs,longtable,array,enumitem}
\usepackage{xcolor}
\usepackage[colorlinks=true,linkcolor=blue!40!black,citecolor=blue!40!black,urlcolor=blue!40!black]{hyperref}
\newtheorem{theorem}{Theorem}[section]
\newtheorem{lemma}[theorem]{Lemma}
\newtheorem{proposition}[theorem]{Proposition}
\theoremstyle{definition}\newtheorem{definition}[theorem]{Definition}
\theoremstyle{remark}\newtheorem{remark}[theorem]{Remark}
\newcommand{\Z}{\mathbb Z}\newcommand{\R}{\mathbb R}
\newcommand{\PP}{\mathbb P}\newcommand{\EE}{\mathbb E}
\newcommand{\eps}{\varepsilon}\newcommand{\ind}{\mathbf1}
\newcommand{\Ecal}{\mathcal E}\newcommand{\Lcal}{\mathcal L}
\newcommand{\Ccal}{\mathcal C}\newcommand{\Q}{Q}
\DeclareMathOperator{\Var}{Var}\DeclareMathOperator{\TV}{TV}
\DeclareMathOperator{\essinf}{ess\,inf}\DeclareMathOperator{\Dir}{Dir}
\allowdisplaybreaks[2]\raggedbottom
\setlength{\emergencystretch}{2em}
\setlist{itemsep=3pt,topsep=5pt}
\hypersetup{pdftitle={Removing the moment assumption from a stipulated planar no-bigeodesics theorem},pdfauthor={Myeongjin Kang}}
\title{Removing the moment assumption from a stipulated\newline planar no-bigeodesics theorem}
\author{Myeongjin Kang}
\date{October 18, 2026}
\begin{document}
\maketitle
\begin{abstract}
Taking the pinned OpenAI planar no-bigeodesics theorem and its specified proof
inputs as established, we give a deduction for every nonatomic law on finite
nonnegative edge weights, without a moment or direction condition. A dependent
circuit cap preserves exactly the set of geometric bigeodesics and supplies
an integrable metric. The capped law does not have upward finite energy.
We replace the required finite-energy operations by a dominated coupling of
lifted limiting graphs, a transfer of labels and selected rays under
underlying-coordinate modifications, and a changed-vertex barrier for
independent half-plane splices. A cubic wall completes the coalescence surgery.
The final budget and perturbation argument is the stipulated source's argument
with status-preserving probability kernels. The additional class is precisely
the laws with infinite second moment of the minimum of four independent
weights. Classical contour surgery, shape theory, and the source's final
counting argument are credited separately from these replacement steps.
\end{abstract}
\noindent\textbf{Reproducibility.} The mathematical deduction is conditional on the stipulated source identified
in Section~\ref{sec:intro}. The full theorem is not formally verified.
The accompanying \texttt{README.md} and \texttt{build.py} give a fresh-directory
PDF build; Section~\ref{sec:coverage} describes the exact limited Lean support.
Original prose is licensed CC BY 4.0; original supporting code is MIT licensed.
No institutional affiliation is asserted.
\tableofcontents
\section{Scope, stipulated input, and attribution}\label{sec:intro}
Let $\Ecal$ be the unoriented nearest-neighbor edges of $\Z^2$. For positive
weights $w$, write $T_w(x,y)$ for the infimum of the costs of finite walks
from $x$ to $y$, counting every traversal. A ray is an infinite simple path
whose finite segments minimize cost; a bigeodesic is such a path indexed by
$\Z$. All weights in this paper are finite. A nonatomic nonnegative law
assigns probability one to strictly positive weights simultaneously on
the countable edge set.

\begin{theorem}\label{thm:main}
Let $G$ be any nonatomic Borel probability measure on $[0,\infty)$, and let
$(t_e)_{e\in\Ecal}$ be independent with common law $G$. Under the stipulated
source convention below, almost surely there is no bigeodesic for $T_t$.
The assertion holds simultaneously over all bi-infinite simple paths,
without a prescribed direction or any moment assumption.
\end{theorem}

Throughout, \textbf{S} denotes OpenAI's \emph{No bigeodesics in planar
first-passage percolation}, September 24, 2026, at immutable revision
\texttt{fd4aeeb2ee4fc729c18d98444fed42fd0529eeeb} of the OpenAI math
repository~\cite{S}. Its theorem and its proof inputs under their stated
hypotheses are stipulated mathematical inputs to this deduction. We do not
audit that upstream theorem and do not apply its iid conclusion to a
dependent field. We use its actual geometric, measurable-selection, and
counting proof operations only after supplying the interfaces below.
\textbf{AH} means Ahlberg--Hoffman, arXiv:1609.02447v2~\cite{AH};
\textbf{DH} means Damron--Hanson, arXiv:1209.3036v1~\cite{DH}.
Their primary results are invoked with the hypothesis distinctions in
Section~\ref{sec:uses}.

S, Theorem 1.1, already covers \emph{every} nonatomic $G$ for which
\begin{equation}\label{eq:paid}
 \EE\big[\min(X_1,X_2,X_3,X_4)^2\big]<\infty,
 \qquad X_i\text{ independent with law }G.
\end{equation}
This includes bounded laws, unbounded finite-mean laws, and some
infinite-mean laws. The additional class in Theorem~\ref{thm:main} is exactly
the class where the expectation in \eqref{eq:paid} is infinite. We may
therefore assume below that $G$ is unbounded. No positive moment is then
assumed. The full stated scope agrees with the finite-valued formulation of
Damron--Hanson's Conjecture 1 in arXiv:1512.00804v2~\cite{DHb}.
This comparison identifies the scope; the deduction here remains conditional
on the stipulated input S.

\paragraph{Why an integrable surrogate is needed.}
The minimum moment in \eqref{eq:paid} is necessary for the usual almost-sure
uniform shape theorem for the original metric. We use the account of
Kesten's necessity and Cox--Durrett/Kesten contour regularization in
Cerf--Th\'eret~\cite[Section 1.2]{CT}; the original Cox--Durrett and Kesten
bodies have not been read for this manuscript. Thus no uniform original
shape or original Busemann linearity is asserted for the additional class.
The cap below has metric regularity, but fails upward finite energy
(Proposition~\ref{prop:energy}). Neither metric is an immediate black-box
input to all the structural results needed by S.

\paragraph{What is new in this deduction.}
The replacement steps are the dominated lifted limiting-graph coupling
(D, Section~\ref{sec:reference}), label and selector transfer under the
underlying-coordinate coupling (S, Section~\ref{sec:labels}), the changed-vertex
half-plane transfer (P, Section~\ref{sec:splice}), and the short cubic wall
(W, Section~\ref{sec:wall}). The circuit cap, exact bigeodesic equivalence,
shape estimates, same-status lowering, two-status density formulas, and S's
budget/fork/tilt argument receive no separate novelty claim.

Open-circuit surgery and its moment-free use have classical antecedents:
Auffinger--Damron--Hanson~\cite[Appendix C, Lemma C.5 and (26)]{ADH} use
open half-circuits to bypass costly geodesic portions; Cox--Durrett/Kesten
contours regularize passage times as described in~\cite{CT}. The direct
wall antecedent is AH, Proposition 9.10, Case 1. The stationary shape theorem
through Boivin is credited through AH Section 3.1 and DH's A2 discussion;
Boivin's original proof is not claimed read. Section~\ref{sec:metric}
retains an elementary shape derivation for completeness, not as a new result.

\paragraph{Logical order.}
First build the cap and its metric estimates. Construct coalescing reference
families by the modified DH argument. Reuse AH's order/label proofs, replacing
the downward step and then the upward selector step. This yields the all-ray
structure. Establish the independent-splice interface, then reuse S's
budget and divergent-toll contradiction with status-preserving kernels.
Every geometric bigeodesic is unchanged by the cap, so the conclusion returns
to the original metric. Earlier clipping, linear-cap, retention, and penalized-path
screens are historical material only and are not premises of this proof.

\section{The cubic circuit cap}\label{sec:cap}
Put $a=\essinf G$. Choose $M\ge1$ so large that $q=G((M,\infty))<1/1024$.
Choose $a<a_o<b_o<M$ with $G([0,a_o])>0$ and $G([0,b_o))<1/1024$.
Nonatomicity at $a$ permits this choice even for singular continuous $G$.
Unbounded support permits $M<a_c<b_c<\infty$ with
$G((M,a_c])>0$ and $G([b_c,\infty))>0$. Set
\begin{align*}
 L_o&=[0,a_o],&H_o&=[b_o,M],&r_o&=G([0,M]),\\
 L_c&=(M,a_c],&H_c&=[b_c,\infty),&r_c&=G((M,\infty)),\\
 p_{a,i}&=G(L_i),&p_{b,i}&=G(H_i),&
 g&=\min(b_o-a_o,b_c-a_c)>0,\\
 p_*&=\min_{i\in\{o,c\}}p_{a,i}/r_i>0.
\end{align*}
All four band masses are positive. Call an original edge \emph{usable} when
$t_e\in H_o\cup H_c$. Its bad probability $\epsilon<1/512$ satisfies
$8\sqrt\epsilon<1$ and $6\sqrt\epsilon<1/2$. These bands remain fixed.

Let $V_e$ be independent uniform variables on $(0,1)$, independent of $t$;
write $X=(t,V)$. An edge is open when $t_e\le M$. Let $L_e$ be the minimum
length of an open simple circuit strictly enclosing both endpoints of $e$,
and define
\begin{equation}\label{eq:cap}
 K_e=M(L_e^3+V_e),\qquad s_e=\min(t_e,K_e),\qquad T=T_s.
\end{equation}
The definition uses the minimum length only, not a choice among tied circuits.

\begin{lemma}\label{lem:capfacts}
Almost surely all $L_e$ exist, $L_e\ge10$, and
$\PP(L_e>r)\le C\exp(-c\sqrt r)$. The field $s$ is a square-symmetric
factor of iid, mixing under every nonzero lattice translation, has every
finite moment, and has unique finite geodesics. Every open edge is unchanged.
\end{lemma}
\begin{proof}
Failure of an open circuit in a suitable annulus of outer radius $R$ forces
a closed dual crossing of length at least $cR$. Erasing loops and counting
at most three choices after the first edge bounds this probability by
$CR\sum_{j\ge cR}(3q)^j$. A simple circuit in the outer box has $O(R^2)$
edges. This gives the stated tail and, by countability, existence everywhere.
The same argument surrounds every finite box by open circuits arbitrarily
far away. A lattice circuit strictly enclosing a horizontal edge has
horizontal span at least three and vertical span at least two; hence its
length is at least ten. The vertical case is identical. In particular
$K_e\ge2ML_e>M$, so open edges are unchanged and
$0<s_e\le M(L_e^3+1)$ has all moments.

Conditional on all open indicators, the original magnitudes have independent
nonatomic conditional laws, the $V_e$ remain independent uniforms, and the
$L_e$ are fixed. Thus the $s_e$ are conditionally independent and nonatomic.
For two distinct simple paths, condition on an edge in their symmetric
difference to rule out equal sums; countability treats all finite paths.
One must not freeze that edge's $V_e$ in this no-tie argument. Existence of
finite minimizers is established by the coercivity argument in
Lemma~\ref{lem:metric}. Covariance and symmetries of the minimum-length
construction give the factor and mixing assertions.
\end{proof}

Choose any measurable minimum circuit $C_e$ for estimates, and put
\[E_x=\{e:x\text{ is strictly inside }C_e\}.\] The tail above and the fact
that $x$ inside $C_e$ implies $\operatorname{dist}(x,e)\le C L_e$ show
that $E_x$ is finite, simultaneously for every $x$.

\begin{proposition}[Exact preservation]\label{prop:equivalence}
On one event, the sets of geometric $t$- and $s$-bigeodesics are exactly
equal. Every edge of each such path is unchanged. Every $s$-ray from a fixed
root $x$ is uncapped outside the finite set $E_x$, simultaneously over all
rays from that root.
\end{proposition}
\begin{proof}
Call $e$ changed if $s_e<t_e$, so $s_e=K_e$. A bigeodesic in either metric
using $e$ must enter and leave $C_e$. Its segment between contacts bracketing
$e$ costs at least $K_e$, whereas a shorter boundary arc costs at most
$ML_e/2$. Thus neither kind of bigeodesic uses a changed edge.

For the nontrivial implication, let $\gamma$ be a $t$-bigeodesic and let
$\pi$ be a finite $s$-walk between $a,b\in\gamma$. Eliminate a changed
occurrence $e$ of $\pi$. If both endpoints lie outside or on $C_e$, replace
the bracketed crossing by a boundary arc. If just $a$ lies inside, replace
the prefix through $e$ by a segment of $\gamma$ from $a$ to a boundary
contact and then a boundary arc to a later contact of $\pi$; reverse the
construction if just $b$ is inside. The part of $\gamma$ between its
outermost contacts with $C_e$ costs at most $ML_e/2$, by $t$-minimality,
so the inserted portion costs at most $ML_e$. If both endpoints are inside,
replace all of $\pi$ by $\gamma[a,b]$, contained in that same outermost-contact
interval and costing at most $ML_e/2$.

Every removed portion contains an edge of $s$-cost $K_e\ge2ML_e$, and every
inserted edge lies on the unchanged circuit or the unchanged $\gamma$.
The number of changed occurrences strictly decreases, with no recursive
creation of changed edges. After finitely many steps the walk is entirely
unchanged and costs no more than the original $s$-walk. Its $t$-cost equals
its $s$-cost, so $t$-minimality of $\gamma$ proves $s$-minimality.
Conversely an unchanged $s$-geodesic is $t$-minimal because $s\le t$.
The proof is pathwise and takes no union over an uncountable path space.

If a ray from $x$ uses changed $e\notin E_x$, it begins outside or on $C_e$
and eventually exits that finite region. The same boundary shortcut is a
contradiction. This proves the rooted-ray statement without asserting any
reattachment of a preserved tail to the same root in the other metric.
\end{proof}

\begin{lemma}[Finite influence and monotonicity]\label{lem:influence}
On one original full event, any finite modification of the positive finite
original weights, with $V$ fixed, changes only finitely many caps. If the
original weights increase, then $s'\ge s$. Under such an upward modification,
an uncapped path avoiding the modified original set is unchanged.
\end{lemma}
\begin{proof}
Tail summability gives $L_e/(1+|e|)\to0$ uniformly over distant edges.
Circuits around every finite box ensure that protecting circuits still exist
after any finite modification $F$. For distant $e$, every circuit enclosing
$e$ and meeting $F$ is longer than $L_e$. An old minimizer avoids $F$, so
$L'_e\le L_e$; a new minimizer of that length also avoids $F$, so it was
old-admissible and $L_e\le L'_e$. Thus equality holds far away.
Raising originals only removes open circuits, giving $L'_e\ge L_e$,
$K'_e\ge K_e$, and $s'_e\ge s_e$. In this upward case an uncapped edge
whose original weight is not modified still has $s'_e=t_e=s_e$.
\end{proof}

\begin{proposition}[Failure of upward finite energy]\label{prop:energy}
For unbounded $G$, the cap law fails DH's upward finite-energy condition
and AH A2(v).
\end{proposition}
\begin{proof}
For $f\ne e$, its status is visible from $s_f$: $s_f\le M$ exactly when
$t_f\le M$. A circuit strictly enclosing both endpoints of $e$ cannot
contain $e$, so $L_e$ is determined by the outside capped field $s_{-e}$.
Consequently $s_e\le M(L_e^3+1)$ conditionally on $s_{-e}$. The event
$L_e=10$ has positive probability and is outside-measurable. On this event
the conditional probability of $s_e\ge\lambda:=1001M+1$ is zero.
Unconditionally that probability is positive: a finite closed pattern
can exclude all short surrounding circuits, an open outer circuit has
positive probability, and independently $t_e>\lambda$ has positive
probability. Take the excluded length sufficiently large that its cap
exceeds $\lambda$. Thus the positive unconditional tail is not conditionally
available everywhere.
\end{proof}
No iid law for $s$ or restoration of its conditional finite energy is used
later. Optional endpoint corrections $T_t-T_s\le B(x)+B(y)$ with
$B(x)=\sum_{e\in E_x}t_e<\infty$ do not supply uniform original shape and
are not needed in the essential chain.

\section{Metric and cocycle estimates}\label{sec:metric}
The cap meets the stationary A2 shape framework discussed in DH and AH,
which credits Boivin. The following elementary derivation is included as
an alternative and to state the precise uniform estimates used below.

\begin{lemma}\label{lem:metric}
Finite geodesics exist uniquely, and there is a deterministic nondegenerate
square-symmetric norm $\mu$ such that, almost surely,
\begin{equation}\label{eq:extended}
 \sup_{x,y\in Q_R}\frac{|T(x,y)-\mu(y-x)|}{R}\longrightarrow0,
 \qquad Q_R=[-R,R]^2\cap\Z^2.
\end{equation}
For every $\eta>0$, eventually uniformly for $x,y\in Q_R$,
$T(x,y)\le C(|x-y|_1+\eta R)$, with $C$ independent of $\eta$.
There is a fixed $L$ such that all geodesics between vertices of $Q_{2n}$
have length at most $Ln$ with probability tending to one.
\end{lemma}
\begin{proof}
Search for protecting circuits of length at most $r$ locally, using their
minimum length if present and $r$ otherwise. The resulting cap $s_e^{(r)}$
depends on a radius $O(r)$ neighborhood, is at most $M(r^3+1)$, satisfies
$s_e^{(r)}\le s_e$, and agrees with $s_e$ unless $L_e>r$.
Along a coordinate segment, color edges by their position modulo a number
between twice the dependence radius and $Cr$. Each color contains independent
variables. Hoeffding and a union over the colors give, for length $\ell\ge Cr$,
\[
 \PP\left(\sum_{e\text{ in segment}}s_e^{(r)}>(\EE s_e+1)\ell\right)
 \le Cr\exp(-c\ell/r^7).
\]
Take $r=\lceil\log(n+2)^4\rceil$. A union over $O(n^3)$ coordinate segments
of length at least $\eta n$, together with $O(n^2)$ coupling failures, has
summable probability
\[
 Cn^3r e^{-c_\eta n/r^7}+Cn^2e^{-c\sqrt r}.
\]
Extend shorter segments to length $\eta n$ in a fixed larger box and use
positivity. Two coordinate segments give the asserted local continuity.

Choose $d\in(0,M)$ with $\PP(t_e<d)$ so small that
$\beta=6\sqrt{\PP(t_e<d)}<1/2$. Since $s_e\ge\min(t_e,M)$, a simple
path of length $\ell$ with cost below $d\ell/2$ has at least half its
independent originals below $d$. Counting paths and bad subsets bounds the
probability of any such path from $O(n^2)$ box vertices, of any length
$\ell\ge n$, by $Cn^2\sum_{\ell\ge n}\beta^\ell$. This is summable.
It supplies coercivity from each fixed vertex and hence existence of finite
minimizers. Coordinate routes have cost at most $C_1n$ on the high-probability
event above, so every box-to-box minimizer has length at most
$\max(n,2C_1n/d)$. Paths leaving the box are included in the lower-cost count.

Kingman's theorem, integrability and ergodicity of each nonzero directional
translation give deterministic radial limits on integer rays. Triangle
inequalities and the lower-cost estimate extend them to a nondegenerate norm
$\mu$. A finite rational net and local continuity give uniform radial shape.
For the shifted version fix $\delta,\eta>0$ and choose one deterministic $N$
so that the event
\[
 G_N(v)=\{|T(v,w)-\mu(w-v)|\le\delta|w-v|
                 \text{ for all }|w-v|\ge N\}
\]
has bad probability less than $c\eta^2$. The spatial ergodic theorem for
this fixed event ensures fewer bad centers in $Q_{2R}$ than the number
of vertices in any $\eta R$-ball. Every $x\in Q_R$ is therefore within
$\eta R$ of a good center $v$. Compare $T(x,y)$ with $T(v,y)$ using local
continuity and norm continuity. If $|y-v|<N$, use local continuity directly.
The error is at most $C\eta R+C'\delta R+O(N)$. Countable decreasing choices
of $\eta,\delta$ prove \eqref{eq:extended}. No radius-dependent event was
fed to the ergodic theorem.
\end{proof}

\begin{lemma}[Stationary extensions]\label{lem:cocycle}
Let a translation-invariant extension of the cap law carry an additive
covariant cocycle $B$ with $|B(x,y)|\le T(x,y)$. The extension need not be
ergodic. There is a possibly random, globally translation-invariant linear
functional $\rho$ such that
\[
 \sup_{x\in Q_R}|B(0,x)-\rho(x)|/R\longrightarrow0.
\]
Consequently for every $\eps>0$ there is finite $C_\eps$ with
$|B(x,y)-\rho(y-x)|\le\eps(|x|+|y|)+C_\eps$ for all $x,y$.
\end{lemma}
\begin{proof}
Two-sided Birkhoff on each axis gives $B(0,ke_i)/k\to\rho_i$, with the
same limit in both signs. For a fixed shift $w$, additivity bounds the
difference of the shifted and unshifted increments by
$T(0,w)+T(ne_i,ne_i+w)$. Stationarity and integrability give
\[
 \sum_{n\ge1}\PP(T(ne_i,ne_i+w)>\eps n)
 =\sum_{n\ge1}\PP(T(0,w)>\eps n)<\infty.
\]
Thus each $\rho_i$ is invariant under every translation, by Borel--Cantelli
and countability. Also $|\rho_i|\le\mu(e_i)$.

For fixed $\alpha>0$ let $Q_N(v)$ assert both axis errors at $v$ are at
most $\alpha|k|$ for every $|k|\ge N$. The conditional bad probability
$d_N=\EE[\ind_{Q_N(0)^c}\mid\mathcal I]$ decreases to zero, where
$\mathcal I$ is the full group-invariant sigma-field. Intersect the spatial
ergodic events for all deterministic $N$. On each component choose a finite
$N$ with $d_N<c\eta^2$. For every $x\in Q_R$ there is
$v\in Q_{\eta R}$ such that both $v$ and $v+x_2e_2$ are good: the two bad
counts in $Q_{3R}$ are smaller than the candidate set. This is a deterministic
count after the ergodic events have been fixed.
Follow the cocycle route
$0\to v\to v+x_2e_2\to v+x\to x$. The middle two long axis increments
have error at most $2\alpha R$; short increments contribute $O(N)$ plus
an arbitrarily small linear remainder by local metric continuity. The end
pieces cost at most $C(\eta+\zeta)R$. Let $R\to\infty$, then the fixed
small parameters decrease to zero. Additivity gives the final extended bound.
No mixing of the enlarged law was inferred from mixing of its cap marginal.
\end{proof}

\section{Reference families by a dominated lifted coupling}\label{sec:reference}
We first distinguish two uses of DH. Its Section 3 averaging construction,
Propositions 5.1--5.2, Theorem 5.3 and Corollary 4.7 apply under A2$'$:
translation ergodicity, square symmetries, a $2+\eps$ moment, bounded norm
shape, and unique passage times. Lemmas~\ref{lem:capfacts}--\ref{lem:metric}
supply these for $s$. DH Section 6, including Lemma 6.2 and the following
triple extraction, has a \emph{standing upward-finite-energy assumption}.
We reuse the proofs of those geometric extraction steps, which use only
supporting-arc geometry, stationarity and the ergodic theorem; we do not
invoke their statements under an unavailable hypothesis. The later
finite-energy surgery is replaced below.

Use DH's exact rounded-line targets
\[
 \widehat L_\alpha=\{z\in\Z^2:(z+[-1/2,1/2)^2)
                  \cap\{x:\rho(x)=\alpha\}\ne\varnothing\}.
\]
No identification of line-target and half-plane-target graph laws is needed.
Coercivity and uniqueness give finite geodesics to each such target. The DH
averages, lifted to include $X$, $s$, graph and cocycle coordinates, have a
stationary subsequential limit $\nu$. Its paths are geodesics, have outdegree
one, merge on intersection, and lie in the supporting arc of the limiting
cocycle slope. The geometric proof of DH Lemma 6.2 and its subsequent
extraction produce, on noncoalescence, a positive-probability triple of
disjoint rays eventually in every translate of one of eight half-planes.
Ergodicity of the lifted law is not needed: its relevant ergodic average
may be random with positive expectation.

\begin{lemma}[D: domination and selected-ray stability]\label{lem:lift}
Fix finite $F\subset\Ecal$ and $a$ with $r=G((a,\infty))>0$.
Resample $t_F$ independently from $G(\cdot\mid t>a)$, keeping everything
else, including $V$, fixed; denote the resulting field by $X'$. There is
a joint limiting coupling with old marginal $\nu$ and new marginal
$\nu'\le r^{-|F|}\nu$. Its underlying pair marginal is exactly this
original-coordinate coupling. On $\max_Ft<a$, any entirely uncapped old
limiting ray avoiding $F$ is the selected new limiting ray from the same root.
\end{lemma}
\begin{proof}
The new original marginal has density
$r^{-|F|}\prod_{e\in F}\ind_{\{t'_e>a\}}\le r^{-|F|}$ relative to the
old iid law. Recompute every finite-line graph and cocycle with the same
$\alpha$ in the two fields. The density bound survives these pushforwards
and averaging over $\alpha$. Joint tightness follows from the fixed
underlying pair law, compact graph coordinates, and cocycle increments
bounded by cap distances; the new distance laws are dominated by a fixed
multiple of the old ones. Take a joint subsequential limit. Domination
passes first for nonnegative bounded continuous functions, then for all
Borel sets by regularity.

Cap coordinates agree with the cap of the retained originals in the limit.
Indeed a shortest circuit and all potentially shorter circuits depend on
finitely many indicators. None of their originals equals $M$ almost surely,
so $L_e$ is locally constant and $s_e$ continuous at almost every original
configuration. Absolute continuity supplies the same fact for the new marginal.

For selected-ray stability fix a root $v$, a finite desired prefix and
$\eps>0$. By finite influence choose $R$ containing $F,v$ so the cap-change
set lies in $Q_R$ with probability at least $1-\eps$. Choose $S\ge R$ so
all old geodesics between vertices of $Q_R$ stay in $Q_S$ with probability
at least $1-\eps$. There are finitely many endpoint pairs; coercivity gives
such an $S$. Both failure bounds are uniform over the approximants because
their underlying marginal is fixed. On this event a path from $v$ cannot
return to $Q_R$ after leaving $Q_S$.

Realize the joint convergence by Skorokhod. Every edge of the asserted old
ray satisfies the strict inequality $t_e<K_e$ almost surely; equality has
probability zero by the conditional nonatomic variables. Its finite prefix
through the first exit from $Q_S$ agrees in late old graph approximants,
and the strict margins and avoidance of $F$ persist. On $\max_Ft<a$ the
change is upward. This prefix is unchanged; on the two localization events
the rest of the finite target path is outside $Q_R$ and unchanged as well.
The whole finite path remains the unique new minimizer to the same target.

If the limiting desired prefixes differed, they would eventually differ
in the discrete approximants. On the asserted limiting event this forces
a localization failure in every sufficiently late approximant. Fatou bounds
the bad limiting probability by $2\eps$. Let $\eps\downarrow0$ and intersect
over roots and prefix lengths. This does not assume almost-sure eventual
localization in the Skorokhod realization.
\end{proof}

\begin{remark}\label{rem:kernels}
The same proof applies to a kernel changing a fixed deterministic finite
set of original coordinates whose
output marginal has a bounded density relative to the underlying product
law, on an event where that kernel increases the originals. The joint
underlying law is fixed across the approximants, the density bound gives
the same domination, and strict uncapping plus the two tight localization
events give the identical stability conclusion. If its upward event is
only measurable, include that event's indicator in the joint lift. Its
discrete convergence makes the event persist in the approximants, where
monotonicity holds; the entire underlying pair-and-indicator marginal is
fixed at every approximation and hence in the limit. In particular,
same-status resampling has output density one, while each finite tilt has
a bounded product density. This extends the explicit reference and label
couplings below to those kernels; it does not assert cap conditional finite
energy or stability under arbitrary uncontrolled modifications. The later
edges and regions are deterministic once their site and radius are fixed;
there are only countably many such choices. A random support is not being
covered by an unstated uniform kernel bound.
\end{remark}

The line convention is compatible with the circuit shortcut: if a rounded
target vertex lies strictly inside a lattice circuit, its half-open unit
cell is inside the polygon. The line enters the polygon and exits at a
boundary edge or vertex, where rounding gives a boundary target vertex.
A cheap boundary arc therefore reaches a valid target. This also justifies
the uncapping argument for the precise DH approximants.

\subsection{W: the cubic wall and its diagonal version}\label{sec:wall}
Suppose the limiting graph is not coalescing. Cut the extracted triple at
last contacts with a sufficiently far common level, after all capped edges.
Countably many choices of roots, level and orientation, followed by
stationarity, give a deterministic positive-probability event of three
wholly uncapped rays rooted on that level and strictly to its positive side
afterward. Selector consistency follows from outdegree one and merging.

For an axis level take roots $(0,j_1),(0,j_2),(0,j_3)$ with
$0<j_1<j_2<j_3<m$, and let $F_m$ consist of the rungs
$(-1,j)(0,j)$, $1\le j\le m-1$. Raise their originals above
$H_m=4Mm^3$. A protecting open circuit enclosing any rung has crossings
of $x=-1/2$ on both sides of it. The intermediate crossings are all closed,
so the vertical span is at least $m$. Hence $L'_e\ge2m$ and
$K'_e\ge8Mm^3>H_m$; the new rung cap is larger than $H_m$.

With probability tending to one, an open circuit strictly surrounding this
wall has at most $C_0m^2$ edges. It avoids the raised set and costs at most
$C_0Mm^2<H_m$ for large $m$. Choose such a deterministic $m$ with positive
joint probability of the triple and cheap circuit, then choose finite
$a\ge H_m$ so that $\max_{F_m}t<a$ still leaves positive probability.
Resample above $a$ as in Lemma~\ref{lem:lift}. The three selected rays persist.

The outer two rays and the level segment joining their roots separate the
plane. A vertex outside the cheap circuit and outside the right-hand corridor
cannot be in the middle class: a path to a sufficiently late middle-ray
vertex would meet an outer ray, merging with a different class, or cross a
rung. In the latter case its segment between two contacts with the cheap
circuit can be bypassed at smaller cost. Thus the middle class lies in the
corridor plus a finite set. It is bounded on the left and meets every fixed
vertical column finitely, since the bounding rays eventually lie beyond
every such column. It has a least vertex in lexicographic $(x,y)$ order.

For a diagonal level use $h=x+y$, $u=x-y$, with roots $(k,-k)$ on $h=0$.
Raise both incoming edges from $h=-1$ into each $(k,-k)$, $|k|\le m$.
Their crossings of $h=-1/2$ have $u=2k-1/2$ and $u=2k+1/2$: consecutive
half-integers from $-2m-1/2$ to $2m+1/2$. A protecting open circuit must
cross on both sides of this blocked interval; its $u$-span exceeds $4m$,
so $L'_e\ge4m$ and $K'_e\ge64Mm^3>H_m$. The same $H_m=4Mm^3$ pays the
$O(Mm^2)$ outer circuit. Entry into the corridor through the interior of
the level segment uses one of the raised incoming edges; entry at an outer
root meets an outer ray. Lexicographic $(h,u)$ order gives the least vertex.
This parity calculation is explicit; a $45$-degree rotation is not being
treated as a square-lattice symmetry. Actual square rotations and reflections
cover the remaining six normals.

Under the stationary old graph law no class can have a least vertex in a
fixed translation-invariant order. Send one unit from each vertex to its
class's least vertex if one exists. Outgoing mass is at most one, whereas
that least vertex receives infinitely many units from its forward ray.
The mass-transport identity and countability rule this out. Domination
$\nu'\le r^{-|F|}\nu$ transfers the null event to the new law, contradicting
the construction. The limiting graph therefore coalesces.
This is the replacement of DH Section 6 surgery and, later, of AH 9.10
Case 1; no cap finite energy has been invoked. The ceiling condition of
DH Lemma 6.4 is vacuous here because the cap has unbounded essential support.

\subsection{Deterministic reference families and the explicit coupling}
For \emph{each tangent functional} $\rho$, choose a point of $\partial\{\mu\le1\}$
at which $\{\rho=1\}$ is the unique supporting line. This is AH's meaning
of tangent; it is precisely the differentiability hypothesis of DH
Corollary 4.7. That result makes the limiting cocycle slope deterministic
$\rho$. Once the graph coalesces its calibrated cocycle is its actual
Busemann function, by evaluating at a common future vertex.

We reuse the proofs of AH 5.3--5.6, not its finite-energy-dependent existence
statement: conditional quantiles of the ordered ray law give, for almost
every fixed quantile, a measurable translation-covariant coalescing family
of slope $\rho$. The already proved graph coalescence supplies the required
premise. The proof of AH 4.4 identifies any two such families. Its boundary
expectation step (AH 4.6) applies to the cap because coordinate, or diagonal
zigzag, boundary routes have finite mean and boundary translations are
ergodic. S's boundary-comparison appendix is unnecessary.

Fix one good quantile. Uniqueness for every other good fixed quantile,
followed by Fubini, says that the conditional ray law given $s$ is a point
mass. Countability of roots identifies the entire graph as a measurable
family $G_\rho(s)$. This statement is conditional on $s$, not merely on
the richer $X$. Domination transfers it to the modified graph.
Apply Lemma~\ref{lem:lift} at late vertices beyond capped edges and $F$.
The old and new $G_\rho$ rays share tails. Since both graphs are now
functions of their cap fields and the projected pair marginal is the
prescribed law of $(X,X')$, this stability holds under that \emph{explicit
original coupling}, not only along a chosen subsequence. Four square
rotations provide the reference families needed for the AH order construction.

\section{Ordered labels and stability of selected rays}\label{sec:labels}
All uses of AH Sections 6--10 in this section are \emph{proof-level reuse}.
AH's standing A1/A2 assumptions include finite-energy conditions not satisfied
by $s$. Section~\ref{sec:uses} lists each use and its replacement premise;
we do not apply its final structural theorem as a black box.

The reference families, uniqueness and planar order give AH Section 6.1's
cyclic ray order. Sections 6.2--6.4 average source flows over independent
auxiliary Voronoi cells. The product structure is the cap law times those
auxiliary uniforms, not independence of capped edges. The conditional-mean
construction in Sections 6.5--6.7 gives monotone labels consistent under
coalescence. Section 7 uses planar trapping, ergodic density and extended
shape to control noncrossing selectors. These arguments introduce no edge
resampling. The one downward operation occurs in AH 8.6.

\subsection{The replacement for AH 8.6}
For an edge $e$, independently resample $t_e$ conditional on its open/closed
status, leaving all $V$ fixed. This kernel is reversible for the underlying
product law. Every status, hence every $L_f$ and $K_f$, stays fixed; only
$s_e$ can change. If $e$ is usable and uncapped, its cap decreases by more
than $\delta=g/2$ with probability at least $p_*$, by sampling its low band.
The fixed $\delta,p_*$ are chosen before the root and finite witness paths
in AH's bad event.

Let $\mathcal F_n$ be the originals and auxiliaries on edges with both
endpoints in $Q_n$. For annular edges outside this set, the kernel preserves
$\mathcal F_n$ and is reversible conditionally on it. Let
$A=A(v,P_1,P_2,P_3,\delta)$ be AH's finite-witness event, with its witness
paths inside $Q_n$. Its deterministic Lemmas 8.4--8.5 use exactly one
decreased edge: a successful decrement on the relevant extremal ray $g_1$
destroys $A$ and places that edge on the new clockwise extremal. They
therefore apply to this exact single-cap-edge change.

Let $C_1(n)$ be the uniform length event of Lemma~\ref{lem:metric}, and let
$C'_1(n)$ use $3Ln$ in place of $Ln$. Lowering one edge implies
$C_1\Rightarrow C'_1$: a new geodesic using the edge splits into two
old-minimal pieces with endpoints in $Q_{2n}$, so has at most $2Ln+1$
edges. An avoiding path is already old-minimal. Usable-edge counting in
the original iid field gives at least half usable edges on every long
simple path from the box, with probability tending to one. All $v$-rays
are uncapped beyond finite $E_v$. Their crossings of the annulus thus have
at least $\gamma n$ eligible edges, say $\gamma=1/3$, on an event $C_2(n,v)$
of probability tending to one.

Writing $D_e$ for eligibility, conditional reversibility gives
\[
 p_*\PP(A\cap C_1\cap D_e\cap\{e\in g_1\}\mid\mathcal F_n)
 \le\PP(A^c\cap C'_1\cap\{e\in g_1\}\mid\mathcal F_n).
\]
Sum over annular edges. On $C'_1$ the part of $g_1$ between its first and
last visits to the outer box has at most $3Ln$ edges. Consequently
\begin{equation}\label{eq:lowering}
 \PP(A\mid\mathcal F_n)\le
 \PP(C_1^c\mid\mathcal F_n)+\PP(C_2^c\mid\mathcal F_n)
 +\frac{3L}{p_*\gamma}\PP(A^c\mid\mathcal F_n).
\end{equation}
The two error terms tend to zero in probability, whereas martingale
convergence makes the conditional probability of $A$ tend to $\ind_A$.
If $\PP(A)>0$, \eqref{eq:lowering} is impossible on a positive-probability
part of $A$ for large $n$. Thus $\PP(A)=0$. Countability of the finite
witnesses yields AH 8.2 with its original simultaneous quantifiers.

The remaining proofs of Sections 8.2--8.4 now use order, compact limits and
countability to give a deterministic closed label set $\Lcal_*$, the two
extremal selectors at each fixed label, and their continuity properties.
The proof of AH 9.2 constructs a stationary increment/path extension.
Lemma~\ref{lem:cocycle}, rather than an assertion of extension ergodicity,
provides its linearity. Its supporting-arc argument and deterministic
selector direction set yield one of the eight eventual half-planes.

\subsection{S: common prediction sites and symmetric label transfer}
Consider the original upward coupling of Lemma~\ref{lem:lift}. It is
generally not stationary as a joint law, but both marginals are absolutely
continuous with respect to the original law. Its cap difference set is finite.
Reference stability from Section~\ref{sec:reference}, together with the
geometric proof of AH 9.7 and extended metric/Busemann bounds, gives a
translate of a cone of positive deterministic angle in which the two
bounding reference rays agree and the sector contains no changed edge.
For agreement, the old and new reference rays avoid the change set and
share a tail; their prefixes to a common tail vertex have the same costs
in both fields, so finite uniqueness identifies them. This replaces the
specific AH 9.6 conclusion used in the proof of 9.7.

Fix a label and the finite neighboring-selector list used in AH 9.8.
Choose a finite stub length so neighboring unequal labels have different
stubs with probability at least $1-\eta$, then a finite cap-patch radius
whose maximum-likelihood predictions of those stubs are simultaneously
correct with probability at least $1-\eta$, adjusting $\eta$ for the
finite list. Martingale convergence supplies this prediction property;
fixed tie-breaking makes the predictors translation-covariant.
Ergodicity gives bad-site density at most $\eta$ in each marginal on a
full event. Absolute continuity transfers that full event to both fields.
Their joint bad-site union therefore has upper density at most $2\eta$,
even though the coupling is not stationary. Choose $2\eta$ below the cone's
density, and discard the finitely many patches meeting the changed set.
Some cone site is good in both fields and has identical patches, hence
identical actual selector stubs.

The planar comparison in the proof of AH 9.8 now traps new extremal labels
between their prescribed old neighbors. New realized labels belong to
$\Lcal_*$ by absolute continuity of the new marginal. If $\alpha$ is
approached from below, use a countable sequence of lower labels. If it has
an immediate predecessor $\beta$, the strict comparison ``new label
$>\beta$'' and membership in $\Lcal_*$ imply new label $\ge\alpha$.
Use the symmetric successor convention above; reference-equality cases
are handled directly. This avoids a nonexistent approximating sequence
at an isolated label. Monotone squeezing transfers every intervening
$\alpha$-ray, not only the two extremals.

The premises of this label comparison are symmetric: both marginals are
absolutely continuous, the cap difference is finite, and corresponding
reference rays coalesce. Apply it in both orientations. The reverse-label
statement does not require a reverse upward resampling kernel.

\begin{lemma}[Finite-change tails]\label{lem:tails}
If two positive coercive metrics $w,w'$ agree off finite $F_0$, every
$w$-ray has a $w'$-geodesic tail, and conversely. The cutoff is uniform
over rays from a fixed root.
\end{lemma}
\begin{proof}
If $F_0$ is empty the assertion is immediate. Otherwise let $V_0$ contain the endpoints of $F_0$, let
$\Delta=\sum_{F_0}|w'_e-w_e|$, and put
$C=\max_{v\in V_0}T_w(x,v)$. A simple competitor from a late $a$ to a later
$b$ on an $x$-ray that meets $V_0$ has new cost at least
\[
 T_w(x,a)+T_w(x,b)-2C-\Delta
 =T_w(a,b)+2T_w(x,a)-2C-\Delta.
\]
Use the first and last contacts and triangle inequalities for this bound.
Once $T_w(x,a)>C+\Delta/2$ and the tail avoids $V_0$, it cannot improve the
unchanged ray segment. Avoiding competitors have identical costs. Positive
weights permit loop erasure before the one-copy $\Delta$ estimate.
Coercivity supplies the uniform late cutoff; reverse the metrics for the converse.
\end{proof}

\subsection{Entire extremal selectors and the second wall application}
Suppose an old extremal $\alpha$-ray from $v$ is wholly unchanged under
$s'\ge s$. It is a new ray by monotonicity. If the new extremal differed
on the relevant side, Lemma~\ref{lem:tails} and the reverse label comparison
would make its tail an old $\alpha$-ray. The proof of AH 9.9 then says old
finite geodesics from $v$ to late vertices of that tail converge to the old
extremal. Put $v$ and the cap changes in $Q_R$, and take $Q_S$ containing
all old geodesics between vertices of $Q_R$. Once such a finite geodesic
agrees with the old extremal through its first exit from $Q_S$, it is wholly
unchanged: a later return to $Q_R$ is impossible. Finite uniqueness makes
it strictly cheaper than the different new-extremal prefix in $s$, and its
unchanged cost is still strictly cheaper in $s'$, a contradiction.
Thus the entire selected ray is preserved. Prefix convergence alone would
not have sufficed without this no-return argument.

Apply the wall of Section~\ref{sec:wall} to late uncapped extremal triples.
The selector premise is now proved by the preceding paragraph, rather than
left as a historical condition. AH 9.11's mass transport gives coalescence.
The remaining proofs of 9.1 and 9.12--9.13 use deterministic labels,
stationarity and disintegration, without a new energy operation.

\subsection{The all-ray output needed by S}\label{sec:allray}
The proof-level chain in AH Section 10, detailed in Section~\ref{sec:uses},
now supplies a deterministic closed set $\Ccal$ of linear functionals
dominated by $\mu$. On one event every cap ray $p$ has a Busemann function
$B_p$ with
\[
 B_p(0,x)=\rho_p(x)+o(|x|),\qquad\rho_p\in\Ccal,
\]
uniformly at infinity; its limiting directions lie in the contact arc
$\{u:\rho_p(u)=\mu(u)\}$. The uniform all-ray cone bounds come from
countably many finite reference nets, not an uncountable intersection of
fixed-ray assertions. For each \emph{fixed} functional there is a unique
coalescing, backwards-finite family; its null set may depend on that
functional. Backwards finiteness uses the mass-transport proof of AH 4.5.
The presence of every tangent functional uses the explicitly quantified
construction in Section~\ref{sec:reference}. All these statements concern
$s$ and its Busemann functions, not uniform asymptotics for the original $t$.

\section{P: independent splices and the changed-vertex barrier}\label{sec:splice}
S Section 3 uses only square-lattice rotations and reflections for its charts.
Fix one chart with upper labels having contact directions strictly above
the horizontal axis and lower labels strictly below it. These are D4
isometries, so coordinate Lipschitz constants and nearest-neighbor height
counts are unchanged. The diagonal wall coordinates used earlier are a
separate construction. The angular gap may depend on a label; no uniform
gap over the label set is asserted.

S Section 4 transfers exceptional labels by independent half-plane pasting.
Its exact shared-weight premise fails for caps. We now prove its required
transfer conclusion with actual changed vertices in place of the cut.
Let $E_n^+$ contain edges whose two endpoint heights are at least $-n$.
Form $X_n^+$ from $X$ on $E_n^+$ and an independent field elsewhere; form
$X_n^-$ from another independent field on $E_n^+$ and $X$ elsewhere.
Splice both originals and cap auxiliaries. For each fixed $n$, the two outputs
are independent and each has the original product law. Their caps $s_n^+,s_n^-$
are therefore independent with the cap law. Independence across $n$ is not needed.

Let $D_n^+$ be all endpoints of edges where $s$ and $s_n^+$ differ, and
define $D_n^-$ similarly. If an unchanged original edge has protecting
minimum circuits in both fields entirely inside the shared region, the
two minima are mutually admissible and its caps agree. Thus, with $|e|$
the midpoint norm, every changed vertex satisfies
\begin{align}
 z\in D_n^+&\Longrightarrow z_2\le-n+3+
       \max_{e\ni z}\big(L_e(X)+L_e(X_n^+)\big),\label{eq:band}\\
 z\in D_n^-&\Longrightarrow z_2\ge-n-3-
       \max_{e\ni z}\big(L_e(X)+L_e(X_n^-)\big).\nonumber
\end{align}
The constants cover endpoint and midpoint conventions. These sets can be
infinite and need not be curves or connected.

Define
\[
 A_n(X)=\sup_e\frac{L_e(X)}{|e|+n},\qquad
 E_n(T)=\sup_{x,y}\frac{|T(x,y)-\mu(y-x)|}{|x|+|y|+n}.
\]
Uniform sublinear coding radii and \eqref{eq:extended} imply that both tend
to zero almost surely. The same quantities for $X_n^+$ tend to zero in
probability, since each has the same marginal law. Choose one deterministic
subsequence $n_k$ on which their probabilities of exceeding $2^{-k}$ are
summable. Borel--Cantelli gives, on that subsequence, all four errors tending
to zero. Consequently for random $\eps_k\downarrow0$,
\begin{equation}\label{eq:movingband}
 z\in D_{n_k}^+\quad\Longrightarrow\quad
 z_2\le-n_k+\eps_k(|z|+n_k).
\end{equation}
For every fixed $n$, both fields' coding radii also give
$z_2\le o(|z|)+O_n(1)$ on $D_n^+$ and the reflected bound on $D_n^-$,
simultaneously for the countable collection of $n$.

\begin{lemma}[Moving and fixed barriers]\label{lem:splice}
Every finite collection of upper-ray Busemann witnesses in $s$ transfers
to $s_{n_k}^+$ for all sufficiently large $k$, with the same exact labels.
For every fixed $n$, every lower ray has a tail in $s_n^-$ with its exact
original label. These statements hold pathwise for all rays on the
structural full event.
\end{lemma}
\begin{proof}
Fix an upper ray $p$, its functional $\rho$, and a finite set of basepoints
$K$ containing $p_0$. Compactness and the strict upper contact sector give
$\eta,\gamma>0$ such that
\begin{equation}\label{eq:gap}
 \mu(d)-\rho(d)\ge\gamma|d|\quad\text{if }d_2\le\eta|d|.
\end{equation}
For large $k$, \eqref{eq:movingband} implies $|z|\ge n_k/2$ for all
$z\in D_{n_k}^+$. Uniformly for $x\in K$, $|z-x|\ge|z|/2$ and
$|x|+|z|+n_k\le8|z-x|$. Also
$(z-x)_2\le3\eps_k|z|+O_K(1)$, so \eqref{eq:gap} applies.
The two metric errors and the Busemann error are small compared with
$|z-x|$. Hence, uniformly on the infinite changed set,
\begin{equation}\label{eq:movingbarrier}
 \min(T,T_n^+)(x,z)-B_p(x,z)
 \ge(\gamma/2)|z-x|\ge\gamma n_k/8.
\end{equation}
The ray itself avoids $D_{n_k}^+$ for large $k$: its eventual strict upper
sector excludes distant vertices by \eqref{eq:movingband}, and the receding
cut excludes its finite prefix.

For a walk from $x$ to $p_j$ meeting $D_{n_k}^+$, let $z$ be its \emph{last}
vertex in that set. Its suffix has exactly the same edge costs in both
fields. Calibration and the metric bound for $B_p$ show that the new walk
costs at least $T_n^+(x,z)+B_p(z,p_j)$, which exceeds
$B_p(x,p_j)$ by the barrier. The same bound holds in the old field. Walks
avoiding the changed set have exactly equal costs. With $x=p_0$ this preserves
all prefixes of $p$. For any fixed $x\in K$ the calibration defect
$T(x,p_j)-B_p(x,p_j)$ tends to zero, so eventually both minimizers avoid
the changed set and
\begin{equation}\label{eq:distanceshare}
 T_n^+(x,p_j)=T(x,p_j).
\end{equation}
This preserves the requested finite Busemann witnesses. The subsequence
was chosen before $p$, so the argument is simultaneous over rays without
new probability intersections.

To identify the exact label fix $n$ and a preserved upper ray. For any
$c>0$, take $x_2\ge c|x|$. The fixed-cut bound gives
$z_2\le\delta|z|+C_\delta$ on $D_n^+$. For small $\delta$ and large
$|x|$, elementary triangle inequalities yield
\[
 |z-x|\ge\kappa(|x|+|z|),\qquad \kappa=c/(c+8).
\]
Indeed $|z-x|<c|x|/4$ would imply $|z|<2|x|$ and
$x_2-z_2>c|x|/2$, a contradiction when $\delta\le c/8$.
Decrease $\delta$ so $\delta/\kappa<\eta/2$. The displacement lies in
the cone of \eqref{eq:gap}; the fixed-metric and Busemann errors then give
\[
 \min(T,T_n^+)(x,z)-B_p(x,z)
 \ge(\gamma\kappa/4)(|x|+|z|)
\]
for all sufficiently far upper-sector $x$, uniformly over $z\in D_n^+$.
The last-changed-vertex argument gives equality of Busemann values at those
$x$. Use $x=ke_2$ and $x=k(e_1+e_2)$ to identify the linear functional
in the new field as exactly $\rho$. The reflected argument applies for
each fixed $n$ to a sufficiently late tail of every lower ray in $s_n^-$,
using two independent lower directions. Thresholds may depend on the ray
and on $n$, as in S's transfer statement.
\end{proof}

S Section 4's remaining proof uses the countability of exceptional labels,
Borel/analytic witness sets and independence of the two output fields.
Lemma~\ref{lem:splice} supplies exactly its two transfer quantifiers,
including finite witnesses and exact labels; fixed-label coalescence is
available from Section~\ref{sec:allray}. Its four-witness independence
argument therefore gives opposite exceptional coalescence simultaneously.
No exact half-plane agreement of capped weights and no original Busemann
linearity have been assumed.

\section{Status-preserving budgets and the final contradiction}\label{sec:finish}
We spell out the probability-kernel replacements and the precise deterministic
inputs reused from S Sections 3--8. Their charts, label budget, finite-tree
count, posterior selection argument, shared-union detour and final filters
are antecedents, not new claims. All events below live on $(t,V)$ and use
cap distances and cap labels. Remark~\ref{rem:kernels} supplies reference
and label stability for the finite kernels used here on their upward events.

\subsection{Charts, comparison, and the one-edge budget}
S Section 3's deterministic chart construction gives countably many compact
label charts $I$, each with a horizontal order parameter $r(\rho)$, finite
width $W_I=r(\rho_r)-r(\rho_\ell)$, and slope $\sigma>0$. The arm cones are
\[
 Q^\pm(k,H)=\{z:\ \pm(z_2-k_2)\ge\sigma|z_1-k_1|-H\}.
\]
Choose nested finite deterministic endpoint sets $D_j$ dense in chart
order and containing the endpoints of every complementary gap. Consecutive
endpoints define the bins $\Pi_j$. Their widths sum to $W_I$, and distinct
nonendpoint labels eventually separate. Fixed-label backwards finiteness
excludes bigeodesic end labels at the countable union of those endpoints.

Here is the deterministic comparison from S Section 5. Suppose a finite
increase near a site $k$ adds total cost at least $b$ on finite initial
portions of a bigeodesic whose upper label lies in a bin $(\alpha,\beta)$.
Suppose its arms have cone offset $H_p$, and the base upper $\alpha$ and
lower $-\beta$ comparator rays from $y$ have offset $H_c$. If the change
and gain portions lie in the radius-$R$ box about $k$, then for
\[
 y_1-k_1>S=\max\left\{R+(R+H_c)/\sigma,(H_p+H_c)/\sigma\right\}
\]
the comparator rays avoid the changed box, remain geodesics with their
labels, and the chart crossing argument gives
\begin{equation}\label{eq:budget}
 b\le B_\alpha(y,k)+B_{-\beta}(y,k).
\end{equation}
This uses finite uniqueness, cone crossing and calibration. Successive
comparisons telescope, leaving only the horizontal label width. The cap
has every one of these premises by Sections~\ref{sec:labels}--\ref{sec:splice}.

At a grid edge $e$, independently replace the original $t_e$ by a sample
from its status-conditional law. Each individual output has exactly the
original iid law, with the same independent $V$. All statuses and thresholds
stay fixed, so only that cap edge changes. Conditional on the \emph{entire}
individual output and $V$, its hidden old value is in the low band of its
status $i$ with probability $p_{a,i}/r_i\ge p_*$. On a usable output
bigeodesic the edge is wholly uncapped by Proposition~\ref{prop:equivalence}.
Mark when its hidden old value is low. The old value is below $a_i<b_i$
and the common threshold exceeds the output value, so the cap increase
is at least $g$ and no other cap changes.

The source's base-comparator regularity check uses both endpoints of $e$:
a base ray either avoids it or uses it first, in which case its tail starts
at the other endpoint. Monotonicity and label transfer give the required
base cone offset from output regularity. Thus its spacing
$\ell(H)\le C_0(H+1)$ and \eqref{eq:budget} apply. If
$A_e^{\alpha,\beta}(H)$ denotes a usable, $H$-regular output bigeodesic
event with label in that bin, the mark probability and budget imply
\begin{equation}\label{eq:edgecount}
 \PP(A_e^{\alpha,\beta}(H))
 \le\frac{\ell(H)\,[r(\beta)-r(\alpha)]}{g p_*},
 \qquad \EE N_e(H)\le C(H+1).
\end{equation}
The latter follows by partition refinement and monotone convergence; $N_e(H)$
counts distinct usable regular edge labels. Independence of marks is not needed.

S Section 6's rooted finite-tree leaf count credits at least $nF_n/4$
distinct usable edge-label pairs to $F_n$ forks. Its usable-path estimate
is valid here in the \emph{original} iid field: a path of length $j$ has
at least half bad edges with probability at most $2^j\epsilon^{j/2}$,
and the path count contributes at most $4^j$, giving $(8\sqrt\epsilon)^j$.
Cap bigeodesics are wholly unchanged. The simultaneous cap cone bound
and \eqref{eq:edgecount} therefore provide exactly its expectation bound.
First let $n\to\infty$ at fixed cone error and then let that error vanish.
There are no forks. With opposite exceptional coalescence from
Section~\ref{sec:splice}, the source's deterministic conclusion identifies
the same two Busemann functions for all bigeodesics of a common random
label, simultaneously. This is not an intersection of fixed-label null events.

\subsection{Finite raises, likelihoods, and posterior independence}
Assume bigeodesics have positive probability. The source's countable chart
and cone-offset reduction gives fixed $I,\sigma,H_0$ and $p>0$ for a
bigeodesic through the origin with arms in $Q^\pm(0,H_0)$ and nonendpoint
label in $I$. Shrink this event, if necessary, by fixing a usable-prefix
threshold $n_0$; the path count above makes the discarded probability
arbitrarily small. Retain the notation $p$ for the resulting positive
probability. These choices are made before the requested toll $m$.

Let $h(e)=\max(|w_2|,|w'_2|)$ for the endpoints of $e$, and put
\[
 \delta_h=\frac{\eta}{(h+2)\log(h+2)},\qquad
 A_k(R)=\left\{e:h(e)\le R,\ 
 \max_{w\in e}|w_1-k_1|\le1+\frac{h(e)+H_0}{\sigma}\right\}.
\]
Take independent trigger uniforms and samples $Z_{e,i}\sim G(\cdot\mid H_i)$,
independent of $X$ and across edges. If $t_e\in L_i$, raise it to $Z_{e,i}$
with probability $\delta_h/p_{a,i}$; otherwise keep it. Choose $\eta$ so
$\delta_0\le\min_i p_{a,i}/2$. Use the same sample for every region
containing $e$, applying it only once in overlaps. All statuses and all
$V$ remain fixed, so all $K_e$ remain fixed. Cap changes are confined
exactly to the raised edges of the specified finite region.

The one-edge output density relative to $G$ is
\begin{align}
 \lambda_h(t)&=1+\sum_{i=o,c}\left[-\frac{\delta_h}{p_{a,i}}\ind_{L_i}(t)
                     +\frac{\delta_h}{p_{b,i}}\ind_{H_i}(t)\right],\label{eq:density}\\
 \int\lambda_h\,dG&=1,\qquad
 \int\lambda_h^2\,dG=1+\delta_h^2\sum_{i=o,c}
                       (p_{a,i}^{-1}+p_{b,i}^{-1}).\nonumber
\end{align}
There is mass $\delta_h$ moved in each status, not just in one status.
The number of region edges at height $h$ is at most
$(12+8/\sigma)(h+H_0+1)$. Independence gives a product likelihood $L$ with
\[
 \EE L^2\le\exp(C\eta^2),\qquad
 \TV(\mathcal L(Y^{k,R}),\mathcal L(X))
 \le\tfrac12\sqrt{e^{C\eta^2}-1},
\]
uniformly in $k,R$, because
$\sum_h(h+H_0+1)/((h+2)^2\log^2(h+2))<\infty$.
Choose $\eta$ so this is at most $p/8$. Finite union fields are also
absolutely continuous. Cap pushforward cannot increase total variation.

Conditional on the whole original output $Y$ and all cap auxiliaries $V$,
the hidden success indicators are independent, with
\begin{equation}\label{eq:posterior}
 \theta_e=\sum_{i=o,c}\ind_{H_i}(Y_e)
                  \frac{\delta_{h(e)}}{p_{b,i}+\delta_{h(e)}}.
\end{equation}
Indeed the joint success/output measure in $H_i$ is
$(\delta_h/p_{b,i})G(dt)$; divide by \eqref{eq:density}. The product
conditional law extends from finite output cylinders to the full output
sigma-field. Circuit lengths are already functions of the output statuses,
so they introduce no additional conditioning. A witnessing cap bigeodesic
is selected from $(Y,V)$ alone. On a successful usable edge of it,
$Y_e=s_e\le K_e$ and the old original is at most $a_i<b_i\le Y_e$;
the old cap is the old original too. Thus the cap gain is at least $g$.

For every prefix length at least $n_0$, at least half its edges are usable.
The $j$th upper-arm edge has $h(e)\le j$, so Abel summation and
\eqref{eq:posterior} give a conditional expected number of successes
bounded below by a positive constant times
$\sum_{j=n_0}^R((j+2)\log(j+2))^{-1}$, which diverges. Conditional
variance is at most the mean. Fix a near radius $R_0$ with mean at least
eight. For any toll $m$, choose $R(m)$ with mean at least
$\max(2\lceil m/g\rceil,32)$. Chebyshev bounds give, jointly, a near
strict raise and total gain at least $m$ with conditional probability at
least $3/4$. The source's marking calculation then gives probability at
least $p/4$. The horizontal spacing
\[
 \ell>R_0+1+(R_0+H_0)/\sigma
\]
is fixed independently of $m$. Posterior independence is not asserted
after the later base or neighbor filters.

\subsection{Shared unions and telescoping}
A bigeodesic in the individual region at $k$ has every edge of height
$\le R$ inside $A_k(R)$ by its cone width; higher edges are never raised.
Passing to the two-site union therefore leaves all of its edges unchanged,
while competitors only increase. The same holds for the selected path at
the second site. A near raised edge of the second path lies outside the
first region by the fixed spacing, so in the first field it still has its
base cap cost and in the union its cost is strictly larger. The common
Busemann pair for equal labels and the source's vanishing-excess detour
through this strict edge rule out repeated marked labels. This is exactly
the geometric use of the near raise; the large radius pays only the toll.

S Section 8's base-comparator filter loses at most $p/8$, leaving
$q_m\ge p/8$. For each fixed $m$, the comparison separation $S(m)$ is finite.
Refine $\Pi_j$ and discard marks with a same-bin neighbor within $S(m)$.
Distinct labels and finitely many neighbors make this filter disappear
pointwise as $j\to\infty$. Apply \eqref{eq:budget} successively in each
bin; Busemann increments telescope and their linear parts leave the sum
$W_I$ of bin widths. The bound is $q_{m,j}\le\ell W_I/m$. Bounded convergence
therefore yields
\[
             p/8\le q_m\le\ell W_I/m\qquad\text{for every integer }m.
\]
The constants $p,\ell,W_I$ were fixed before $m$, a contradiction for large
$m$. Sites, finite regions, unions, partition endpoints and chosen radii
form countable collections, so all inherited full events can be imposed
simultaneously. Random-label equalities were already pathwise.
There are no cap bigeodesics. Proposition~\ref{prop:equivalence} proves
Theorem~\ref{thm:main} for the remaining laws, and S covers \eqref{eq:paid}.

\section{Proof-use crosswalk and verification limits}\label{sec:coverage}\label{sec:uses}
The following map records proof reuse, not the invocation of a theorem whose
standing finite-energy hypothesis is absent. In each row the original geometric
or order argument is retained only after the indicated replacement has been
supplied. In particular the cap law is never asserted to satisfy AH A2(v) or
A2(vii). All probabilities for coordinate surgery are taken on the underlying
product space, including the independent cap marks.

\begin{longtable}{@{}>{\raggedright\arraybackslash}p{.22\textwidth}>{\raggedright\arraybackslash}p{.72\textwidth}@{}}
\toprule
Original use & Hypotheses and replacement in this deduction \\\midrule\endhead
DH Sections 3--5 & A2$'$ supplies the averaged rounded-line graphs, cocycles,
outdegree, direction and supporting-arc conclusions. Section~\ref{sec:metric}
proves the required cap shape and integrability. The lifted extension is not
assumed ergodic; Lemma~\ref{lem:cocycle} handles its invariant random slope.\\
DH Lemma 6.2 and triple extraction & These occur under the Section 6 standing
upward-finite-energy assumption. Only their planar extraction proofs are reused:
stationarity, directedness and positive ergodic density give three disjoint
rays in one of the eight eventual half-planes. The probability-changing step
is Lemma~\ref{lem:lift} and Section~\ref{sec:wall}, not the DH finite-energy
conclusion. DH 6.4's bounded-ceiling issue is vacuous for this unbounded cap.\\
AH 4.4, 4.6; 5.3--5.6 & Fixed-slope uniqueness uses cap boundary integrability.
Conditional quantile extraction uses the coalescent graph just constructed;
it does not supply coalescence. Point-mass conditional graph laws descend the
joint limiting coupling to the explicit underlying-coordinate pair.\\
AH 6.1--6.7 & The auxiliary partitions and labels use independent auxiliary
randomness, stationarity, shape, finite uniqueness and the four reference
families. Independence of cap edge weights is not used. References are
constructed in Section~\ref{sec:reference}.\\
AH Section 7 & The noncrossing order, extremal rays and limiting labels use
the preceding references, planarity, uniqueness, measurability and shape.
They are reused at proof level in Section~\ref{sec:labels}.\\
AH 8.3--8.6 & The single-edge downward step is replaced by same-status
resampling. It changes exactly one cap coordinate, has a reversible underlying
law, and lowers an uncapped good edge by a fixed positive amount with uniformly
positive probability. The annular good-edge density, eventual uncapping,
conditional estimate and martingale contradiction are given in
Section~\ref{sec:labels}, including equation~\eqref{eq:lowering}.\\
AH 8.1--8.4, remaining conclusions & The resulting deterministic closed label
set and extremal selectors follow by the original order/countability proofs
after the previous row. Membership of new labels in that same set is transferred
by absolute continuity. Isolated predecessors and successors are treated
explicitly, not by nonexistent monotone approximating sequences.\\
AH 9.2--9.4 & Random-extension cocycle linearity follows from
Lemma~\ref{lem:cocycle}; directional ergodicity is not inferred from cap mixing.
The order and supporting-cone consequences then use that linearity and shape.\\
AH 9.5--9.9 & Common local reference stubs have positive density under the
two absolutely continuous marginals. Comparison is symmetric; squeezing
extends it to every ray with the label. Lemma~\ref{lem:tails} supplies the
finite-change tail map. The final no-return argument upgrades preserved tails
to the selected whole ray on the stated upward, unchanged-path event.\\
AH 9.10, 9.11 & The wall operation is Section~\ref{sec:wall}, with cubic cap,
cheap outer circuit, exact diagonal crossings and least-class mass transport.
The required selector stability is established before this reuse. The
subsequent coalescence deduction uses the original planar argument.\\
AH 9.1, 9.12--9.13; Section 10 & These are reused only after the structural
and modification premises above. Countable label approximations, planarity,
coalescence, shape and fixed-functional uniqueness give the all-ray output
stated in Section~\ref{sec:allray}; fixed-label null events are not silently
intersected over an uncountable random set.\\
S Section 04 & Product independence belongs to spliced originals and marks.
Section~\ref{sec:splice} proves moving-witness and fixed-cut transfers across
the changed-cap band, including equality of the late suffix costs and labels.
The remaining exceptional-label countability argument is the source's.\\
S Sections 05--08 & Section~\ref{sec:finish} keeps the budget, fork, grid,
harmonic tilt, neighbor filter and telescoping operations. Two-status kernels
replace iid cap resampling. Their exact densities and posterior probabilities
are equations~\eqref{eq:density}--\eqref{eq:posterior}; statuses and cap radii
remain fixed. No new contribution is claimed for the source's counting argument.\\
\bottomrule
\end{longtable}

Remark~\ref{rem:kernels} explicitly generalizes the fixed-threshold coupling
to fixed finite-support kernels with bounded output density. The underlying
pair (and, when needed, its upward-event indicator) is retained in the lift;
its projected law is identical at every approximation and in the limit.
Selected-path stability is restricted to the upward event. Each later support
is deterministic after the specified integer site and radius are fixed.
Countably many choices may then be intersected; arbitrary state-dependent
supports or unbounded-density kernels are not covered by this statement.

\subsection{Limited deterministic formal support}
The accompanying frozen Lean support has 14 support declarations and 17 control
declarations. Its principal all-walk statement concerns a graph with a
nonnegative changed weight function, a finite set of changed edges, a finite
absolute-change budget, calibrated old ray segments, radial additivity and
escape, and a bound on old root distances to changed-edge endpoints. It proves
that sufficiently late segments minimize over all finite walks. The passage
from edge-simple competitors to walks uses loop erasure and nonnegativity.
This supports the deterministic argument of Lemma~\ref{lem:tails}.

The formal objects do not construct an infinite random geodesic or prove the
probabilistic assumptions supplying calibration, coercivity or finite
influence. They do not formalize the circuit cap, shape, conditional kernels,
Busemann limits, labels, coalescence, splicing, or Theorem~\ref{thm:main}.
The finite controls illustrate hypotheses and counterexamples; their number
is not a count of independently proved probabilistic lemmas.

\subsection{Source-reading boundary}
The author derivation uses the operative source and primary interfaces listed
in the accompanying reading record. Prior proof scrutiny was bounded: in
particular some AH Sections 4, 6, 7 and 10 and DH Sections 3--5 were checked
at statement/interface level rather than through every proof line. This
manuscript makes the required proof-level uses explicit for further scrutiny;
it does not turn such bounded reading into a claim of exhaustive independent
verification. Cox--Durrett, Kesten and Boivin are credited through the actually
read discussions in AH, DH and Cerf--Th\'eret. No original-body reading of
those works is asserted. The public evidence records distinguish author
assistance, independent reports, formal support and compilation checks.

\begin{thebibliography}{99}
\bibitem{S} OpenAI, \emph{No bigeodesics in planar first-passage percolation}, stipulated source 212, repository revision \texttt{fd4aeeb2ee4fc729c18d98444fed42fd0529eeeb}. Exact source paths and blob identities are recorded in the accompanying source receipts.
\bibitem{AH} D. Ahlberg and C. Hoffman, \emph{Random coalescing geodesics in first-passage percolation}, \href{https://arxiv.org/abs/1609.02447v2}{arXiv:1609.02447v2}.
\bibitem{DH} M. Damron and J. Hanson, \emph{Busemann functions and infinite geodesics in two-dimensional first-passage percolation}, \href{https://arxiv.org/abs/1209.3036v1}{arXiv:1209.3036v1}; Communications in Mathematical Physics 325 (2014), 917--963, doi:10.1007/s00220-013-1875-y.
\bibitem{DHb} M. Damron and J. Hanson, \emph{Bigeodesics in first-passage percolation}, \href{https://arxiv.org/abs/1512.00804v2}{arXiv:1512.00804v2}, Conjecture 1.
\bibitem{ADH} A. Auffinger, M. Damron and J. Hanson, \emph{Limiting geodesics for first-passage percolation on subsets of $\Z^2$}, \href{https://arxiv.org/abs/1302.5413v2}{arXiv:1302.5413v2}; Annals of Applied Probability 25 (2015), 373--405, doi:10.1214/13-AAP999. Appendix C, especially Lemma C.5 and (26).
\bibitem{CT} R. Cerf and M. Th\'eret, \emph{Weak shape theorem in first passage percolation with infinite passage times}, \href{https://arxiv.org/abs/1404.4539v2}{arXiv:1404.4539v2}.
\end{thebibliography}

\end{document}
